% ===========================================================================
% CHƯƠNG 1: CƠ SỞ LÝ LUẬN
% ===========================================================================
\chapter{CƠ SỞ LÝ LUẬN}
% Tiêu đề chạy rút gọn: tiêu đề chương đầy đủ dài hơn một dòng đầu trang,
% khiến fancyhdr phải nới \headheight và làm lệch khung chữ của cả chương.
\markboth{Chương 1. Cơ sở lý luận}{}
\label{chap:co-so-ly-luan}

\noindent\textit{Chương 1 hệ thống hóa cơ sở lý thuyết của đề án. Nội dung chương
lần lượt trình bày tổng quan về tín dụng và hoạt động kiểm duyệt tín dụng; tổng
quan về học máy và vai trò của học máy trong đánh giá tín dụng; đặc điểm dữ liệu
và phương pháp xây dựng đặc trưng cho bài toán tín dụng; và cuối cùng là các mô
hình học máy được nghiên cứu trong đề án cùng cơ sở so sánh, lựa chọn mô hình phù
hợp. Đây là nền tảng để Chương 2 xây dựng mô hình và Chương 3 tiến hành thực
nghiệm.}

\section{Tổng quan về tín dụng và kiểm duyệt tín dụng}
\label{sec:tong-quan-tin-dung}

\subsection{Khái niệm tín dụng}
\label{subsec:khai-niem-tin-dung}

Tín dụng là quan hệ chuyển nhượng quyền sử dụng một lượng giá trị -- thường dưới
hình thái tiền tệ -- từ bên cho vay sang bên đi vay trong một khoảng thời gian
xác định, với cam kết hoàn trả cả gốc và lãi khi đến hạn. Ba đặc trưng cơ bản của
quan hệ tín dụng gồm: \textit{tính hoàn trả} (khoản vốn phải được trả lại, khác
với cấp phát hay tài trợ không hoàn lại), \textit{tính thời hạn} (quyền sử dụng
vốn chỉ được chuyển giao trong một thời gian nhất định) và \textit{tính sinh lời}
(bên cho vay nhận được khoản lãi như phần bù cho việc từ bỏ quyền sử dụng vốn và
cho rủi ro phải gánh chịu).

Đặc trưng thứ ba dẫn tới một hệ quả quan trọng về mặt kinh tế: lãi suất không chỉ
phản ánh giá của vốn theo thời gian mà còn phải bao hàm \textit{phần bù rủi ro}.
Một tổ chức tín dụng cho vay với cùng một mức lãi suất đối với mọi khách hàng sẽ
gặp hiện tượng lựa chọn ngược: khách hàng có rủi ro thấp thấy mức lãi suất đó
không hấp dẫn và tìm nguồn vốn khác, còn khách hàng có rủi ro cao sẵn sàng chấp
nhận. Kết quả là danh mục tín dụng dần tập trung vào nhóm rủi ro cao trong khi
lãi suất không được điều chỉnh tương xứng. Đây chính là lý do nền tảng khiến việc
\textit{phân biệt được mức độ rủi ro của từng khách hàng} trở thành năng lực cốt
lõi của hoạt động tín dụng, chứ không đơn thuần là một yêu cầu về tuân thủ.

Trong phạm vi đề án, đối tượng nghiên cứu là tín dụng bán lẻ đối với khách hàng
cá nhân: các khoản vay tiêu dùng, vay mua ô tô, vay thế chấp bất động sản và thấu
chi. Nhóm sản phẩm này có hai đặc điểm khiến chúng đặc biệt phù hợp với phương
pháp mô hình hóa thống kê. Thứ nhất, số lượng hồ sơ lớn nhưng giá trị mỗi khoản
vay tương đối nhỏ, nên chi phí thẩm định thủ công cho từng hồ sơ khó có thể tương
xứng với quy mô khoản vay. Thứ hai, các khoản vay trong cùng một sản phẩm có mức
độ đồng nhất cao, cho phép rút ra quy luật thống kê từ dữ liệu lịch sử -- điều
khó thực hiện với tín dụng doanh nghiệp lớn, nơi mỗi khoản vay gần như là một
trường hợp riêng biệt.

\subsection{Quy trình cấp và kiểm duyệt tín dụng}
\label{subsec:quy-trinh-kiem-duyet}

Quy trình cấp tín dụng tại các tổ chức tín dụng thường được tổ chức thành năm
bước chính:

\begin{enumerate}[label=(\arabic*)]
  \item \textbf{Tiếp nhận hồ sơ.} Khách hàng nộp đề nghị vay vốn kèm các tài liệu
  chứng minh nhân thân, thu nhập, mục đích sử dụng vốn và tài sản bảo đảm (nếu
  có). Hồ sơ có thể được tiếp nhận tại quầy, qua kênh trực tuyến hoặc qua đối tác
  phân phối.

  \item \textbf{Thu thập và kiểm tra thông tin.} Cán bộ tín dụng đối chiếu thông
  tin khách hàng khai báo với dữ liệu nội bộ của ngân hàng (lịch sử quan hệ, số
  dư tài khoản, lịch sử giao dịch) và với dữ liệu bên ngoài, trong đó quan trọng
  nhất là bản tin tín dụng từ Trung tâm Thông tin tín dụng Quốc gia Việt Nam
  (CIC)~\cite{tt15}.

  \item \textbf{Thẩm định và kiểm duyệt.} Đây là bước trọng tâm của quy trình:
  đánh giá năng lực tài chính, khả năng trả nợ, mức độ rủi ro của khách hàng và
  sự phù hợp của khoản vay với chính sách tín dụng hiện hành. Kết quả của bước
  này là đề xuất phê duyệt, từ chối hoặc yêu cầu bổ sung thông tin.

  \item \textbf{Phê duyệt.} Cấp có thẩm quyền -- tùy theo hạn mức được phân cấp
  -- ra quyết định cuối cùng trên cơ sở hồ sơ và đề xuất thẩm định.

  \item \textbf{Giải ngân và giám sát sau vay.} Sau khi hoàn tất thủ tục pháp lý,
  khoản vay được giải ngân và chuyển sang giai đoạn theo dõi việc trả nợ, phân
  loại nợ và trích lập dự phòng rủi ro theo quy định~\cite{tt11}.
\end{enumerate}

Phạm vi của đề án tập trung vào bước thứ ba -- \textbf{kiểm duyệt tín dụng}. Cần
phân biệt rõ hai khái niệm dễ nhầm lẫn: \textit{kiểm duyệt} là quá trình đánh giá
và sàng lọc hồ sơ để hình thành đề xuất có căn cứ, còn \textit{phê duyệt} là hành
vi ra quyết định của cấp có thẩm quyền. Hệ thống được xây dựng trong đề án hỗ trợ
bước kiểm duyệt: nó cung cấp đánh giá rủi ro định lượng kèm lý giải, làm đầu vào
cho quyết định của con người, chứ không thay thế thẩm quyền phê duyệt.

Sự phân biệt này không chỉ mang tính thuật ngữ mà có ý nghĩa thiết kế trực tiếp.
Một mô hình \textit{thay thế} con người buộc phải đưa ra quyết định nhị phân cho
mọi hồ sơ, kể cả những hồ sơ mà bản thân nó không đủ chắc chắn. Một mô hình
\textit{hỗ trợ} có thể -- và nên -- thừa nhận vùng không chắc chắn của mình và
chuyển những hồ sơ đó cho người có thẩm quyền xem xét. Toàn bộ thiết kế ba vùng
quyết định được trình bày ở Chương~\ref{chap:xay-dung-he-thong} bắt nguồn từ nhận
định này.

\subsection{Các yếu tố ảnh hưởng đến quyết định tín dụng}
\label{subsec:yeu-to-anh-huong}

Thực tiễn thẩm định tín dụng từ lâu đã tổng kết các yếu tố cần xem xét thành
khung phân tích \textbf{5C}~\cite{thomas2017}:

\begin{itemize}
  \item \textbf{Character} (Uy tín): thái độ và lịch sử thực hiện nghĩa vụ trả nợ
  của khách hàng. Trong dữ liệu, yếu tố này được phản ánh qua lịch sử tín dụng
  tại CIC: nhóm nợ, số ngày quá hạn, tiền sử nợ đã xử lý rủi ro.

  \item \textbf{Capacity} (Năng lực trả nợ): khả năng tạo ra dòng tiền để trả nợ,
  đo bằng thu nhập, tính ổn định của thu nhập và các chỉ tiêu như tỷ lệ nợ trên
  thu nhập ($\DTI$) hay tỷ lệ khoản trả góp trên thu nhập.

  \item \textbf{Capital} (Vốn tự có): phần vốn khách hàng tự bỏ ra trong phương
  án vay, thể hiện mức độ ``cùng chịu rủi ro''. Với khoản vay có tài sản bảo đảm,
  yếu tố này gắn với tỷ lệ vay trên giá trị tài sản ($\LTV$).

  \item \textbf{Collateral} (Tài sản bảo đảm): nguồn thu hồi nợ thứ cấp khi nguồn
  trả nợ chính không thực hiện được.

  \item \textbf{Conditions} (Điều kiện): bối cảnh kinh tế vĩ mô, đặc điểm ngành
  nghề và mục đích sử dụng vốn.
\end{itemize}

Khung 5C có giá trị định hướng nhưng bộc lộ một hạn chế khi chuyển thành quy tắc
vận hành: nó liệt kê các yếu tố cần xem xét mà không quy định cách \textit{kết
hợp} chúng thành một đánh giá tổng thể. Trong thực tế, mức độ quan trọng tương
đối của từng yếu tố thay đổi theo phân khúc khách hàng và theo thời gian; hơn nữa
các yếu tố tương tác với nhau theo cách phi tuyến. Một khách hàng có $\DTI$ ở mức
0,45 có thể chấp nhận được nếu thu nhập ổn định và lịch sử tín dụng tốt, nhưng
lại rủi ro cao nếu thâm niên công việc ngắn và đang có quan hệ tín dụng với nhiều
tổ chức. Chính khoảng trống này -- xác định trọng số và dạng tương tác giữa các
yếu tố từ dữ liệu thay vì từ phán đoán chủ quan -- là chỗ mà phương pháp học máy
đóng góp giá trị.

\subsection{Rủi ro tín dụng và đánh giá rủi ro tín dụng}
\label{subsec:rui-ro-tin-dung}

\textbf{Rủi ro tín dụng} là khả năng bên đi vay không thực hiện được nghĩa vụ trả
nợ theo cam kết, dẫn tới tổn thất cho bên cho vay. Theo khung Basel
II~\cite{bcbs2006}, tổn thất kỳ vọng của một khoản tín dụng được phân rã thành
tích của ba thành phần:

\begin{equation}
\EL = \PD \times \LGD \times \EAD,
\label{eq:expected-loss}
\end{equation}

trong đó $\PD$ (\eng{Probability of Default}) là xác suất vỡ nợ trong một khoảng
thời gian xác định -- thường là 12 tháng; $\LGD$ (\eng{Loss Given Default}) là tỷ
lệ tổn thất khi vỡ nợ xảy ra; và $\EAD$ (\eng{Exposure at Default}) là dư nợ tại
thời điểm vỡ nợ.

Phân rã này cho thấy vị trí của bài toán mà đề án giải quyết. Trong ba thành
phần, $\PD$ là đại lượng gắn trực tiếp với \textit{đặc điểm của khách hàng} và
được quyết định tại thời điểm cấp tín dụng; $\LGD$ và $\EAD$ phụ thuộc chủ yếu
vào cấu trúc khoản vay, tài sản bảo đảm và quá trình thu hồi nợ sau này. Vì vậy
$\PD$ là đại lượng mà hoạt động kiểm duyệt tín dụng có thể tác động và cũng là
đại lượng mà mô hình học máy có thể ước lượng từ dữ liệu lịch sử. Đề án tập trung
vào việc ước lượng $\PD$; việc mô hình hóa $\LGD$ và $\EAD$ không thuộc phạm vi
nghiên cứu.

\paragraph{Định nghĩa vỡ nợ.} Để ước lượng $\PD$, trước hết phải định nghĩa được
thế nào là ``vỡ nợ''. Thông tư 11/2021/TT-NHNN~\cite{tt11} quy định việc phân
loại nợ thành năm nhóm theo mức độ rủi ro, trong đó nợ thuộc nhóm 3 (Nợ dưới tiêu
chuẩn), nhóm 4 (Nợ nghi ngờ) và nhóm 5 (Nợ có khả năng mất vốn) được xếp vào
\textbf{nợ xấu}. Về mặt số ngày quá hạn, ranh giới của nợ nhóm 3 là 90 ngày. Định
nghĩa này trùng với ngưỡng mà Basel sử dụng cho khái niệm \eng{default}, nhờ đó
kết quả nghiên cứu có thể đối chiếu với các công trình quốc tế. Đề án sử dụng
thống nhất định nghĩa: một hồ sơ được gán nhãn \textit{xấu} nếu phát sinh tình
trạng quá hạn từ 90 ngày trở lên (tương đương nợ nhóm 3--5) trong cửa sổ quan sát
kể từ thời điểm cấp tín dụng.

\paragraph{Ba thế hệ phương pháp đánh giá.} Công tác đánh giá rủi ro tín dụng đã
trải qua ba giai đoạn phát triển. Giai đoạn đầu dựa trên \textit{phán đoán chuyên
gia}: cán bộ tín dụng vận dụng kinh nghiệm để đưa ra nhận định, linh hoạt với
trường hợp đặc thù nhưng thiếu nhất quán và khó mở rộng theo quy mô. Từ thập niên
1960, \textit{thẻ điểm thống kê} (\eng{scorecard}) dựa trên hồi quy logistic ra
đời, quy đổi mỗi đặc điểm khách hàng thành một số điểm; phương pháp này bảo đảm
tính nhất quán, cho phép kiểm định thống kê và dễ giải thích cho cơ quan quản lý.
Từ thập niên 2000, \textit{mô hình học máy} được áp dụng ngày càng rộng rãi nhờ
khả năng nắm bắt quan hệ phi tuyến và tương tác bậc cao mà không cần chỉ định
trước dạng hàm, đổi lại khó giải thích hơn và đòi hỏi kiểm định chặt chẽ hơn.

Ba thế hệ này không thay thế nhau hoàn toàn. Trong thực tiễn hiện nay, các tổ
chức tín dụng thường vận hành đồng thời cả ba: mô hình học máy cung cấp điểm rủi
ro, bộ luật nghiệp vụ bảo đảm tuân thủ chính sách và ngưỡng pháp lý, còn phán
đoán chuyên gia xử lý các hồ sơ nằm ngoài vùng mà mô hình đủ tin cậy. Hệ thống
được xây dựng trong đề án đi theo mô hình kết hợp này.

\subsection{Những hạn chế của phương pháp kiểm duyệt tín dụng truyền thống}
\label{subsec:han-che-truyen-thong}

Phương pháp kiểm duyệt truyền thống -- dựa trên bộ quy tắc nghiệp vụ cố định kết
hợp với phán đoán của cán bộ tín dụng -- bộc lộ bốn hạn chế chính.

\paragraph{Thứ nhất, không biểu diễn được tương tác giữa các yếu tố.} Một bộ quy
tắc có dạng ``nếu $\DTI > 0{,}6$ thì từ chối'' xét từng chỉ tiêu độc lập với các
chỉ tiêu còn lại. Cách làm này ngầm giả định rằng ảnh hưởng của mỗi yếu tố tới
rủi ro là cộng tính và không phụ thuộc giá trị của các yếu tố khác. Giả định đó
thường bị vi phạm trong thực tế: tác động của tỷ lệ $\DTI$ cao lên xác suất vỡ nợ
mạnh hơn nhiều ở nhóm khách hàng có thu nhập không ổn định so với nhóm có thu
nhập ổn định. Về mặt hình thức, bộ quy tắc chỉ tạo ra được các mặt phân chia song
song với trục tọa độ trong không gian đặc trưng, trong khi ranh giới rủi ro thực
tế có dạng phức tạp hơn nhiều.

\paragraph{Thứ hai, ngưỡng được đặt theo phán đoán chứ không theo dữ liệu.} Các
ngưỡng như ``$\DTI$ tối đa 60\%'' hay ``thâm niên công việc tối thiểu 6 tháng''
thường được kế thừa từ thông lệ hoặc từ kinh nghiệm của người xây dựng chính
sách. Chúng hiếm khi được hiệu chỉnh lại theo dữ liệu thực tế của chính tổ chức,
và càng hiếm khi được kiểm định xem mức ngưỡng đó có thực sự tối ưu về mặt đánh
đổi giữa quy mô cấp tín dụng và chất lượng danh mục hay không.

\paragraph{Thứ ba, kết quả đầu ra mang tính nhị phân, không có thang đo mức độ.}
Một bộ quy tắc cho biết hồ sơ ``đạt'' hay ``không đạt'' nhưng không cho biết hồ
sơ đạt ở mức nào. Điều này khiến tổ chức tín dụng không thể thực hiện hai nghiệp
vụ quan trọng: định giá theo rủi ro (khách hàng rủi ro thấp được hưởng lãi suất
ưu đãi hơn) và phân bổ nguồn lực thẩm định (tập trung nhân lực vào nhóm hồ sơ
thực sự cần xem xét kỹ).

\paragraph{Thứ tư, tính nhất quán và tốc độ xử lý bị giới hạn bởi yếu tố con
người.} Cùng một hồ sơ có thể nhận được kết luận khác nhau từ hai cán bộ khác
nhau, hoặc từ cùng một cán bộ ở hai thời điểm khác nhau. Khi khối lượng hồ sơ
tăng, thời gian xử lý kéo dài và áp lực công việc làm gia tăng nguy cơ sai sót.

Bốn hạn chế trên đều có chung một đặc điểm: chúng không phải là vấn đề về năng
lực chuyên môn mà là giới hạn cấu trúc của phương pháp. Do đó, hướng khắc phục
không phải là tăng cường đào tạo hay bổ sung thêm quy tắc, mà là bổ sung một
thành phần có khả năng học quan hệ từ dữ liệu và cho ra kết quả trên thang đo
liên tục. Đó chính là vai trò của mô hình học máy mà đề án xây
dựng.

\section{Tổng quan về học máy trong đánh giá tín dụng}
\label{sec:tong-quan-hoc-may}

\subsection{Khái niệm học máy}
\label{subsec:khai-niem-hoc-may}

Học máy (\eng{machine learning}) là lĩnh vực nghiên cứu các thuật toán cho phép
hệ thống máy tính cải thiện hiệu quả thực hiện một nhiệm vụ thông qua kinh
nghiệm, thay vì thông qua việc được lập trình tường minh từng bước. Định nghĩa
được trích dẫn rộng rãi của Mitchell~\cite{mitchell1997} phát biểu: một chương trình máy tính được
gọi là \textit{học} từ kinh nghiệm $E$ đối với lớp nhiệm vụ $T$ và độ đo hiệu quả
$P$, nếu hiệu quả của nó trên nhiệm vụ $T$, đo bằng $P$, cải thiện theo kinh
nghiệm $E$.

Áp dụng vào bài toán của đề án, ba thành phần này được xác định cụ thể như sau:

\begin{itemize}
  \item \textbf{Nhiệm vụ $T$}: dự báo khả năng một hồ sơ vay sẽ phát sinh nợ xấu
  trong 12 tháng kể từ thời điểm cấp tín dụng.
  \item \textbf{Kinh nghiệm $E$}: tập hồ sơ vay trong quá khứ, mỗi hồ sơ gồm các
  đặc trưng được quan sát tại thời điểm nộp và kết quả trả nợ đã biết.
  \item \textbf{Độ đo hiệu quả $P$}: khả năng phân biệt giữa hồ sơ tốt và hồ sơ
  xấu, đo bằng các chỉ tiêu như $\AUC$, hệ số Gini và thống kê $\KS$.
\end{itemize}

Điểm khác biệt căn bản giữa học máy và lập trình theo quy tắc nằm ở chiều của
quá trình suy luận. Lập trình theo quy tắc đi từ quy luật do con người phát biểu
tới kết quả; học máy đi ngược lại -- từ các cặp (dữ liệu đầu vào, kết quả đã
biết) để suy ra quy luật. Trong bài toán tín dụng, chiều ngược này có ý nghĩa
thực tiễn rõ rệt: thay vì yêu cầu chuyên gia phát biểu chính xác mối quan hệ giữa
hàng trăm đặc trưng và xác suất vỡ nợ (điều vượt quá khả năng của phán đoán chủ
quan), thuật toán rút ra quan hệ đó từ chính dữ liệu lịch sử của tổ chức.

\subsection{Phân loại các phương pháp học máy}
\label{subsec:phan-loai-hoc-may}

Các phương pháp học máy thường được phân loại theo dạng thông tin giám sát có
trong dữ liệu huấn luyện.

\paragraph{Học có giám sát (\eng{supervised learning}).} Dữ liệu huấn luyện gồm
các cặp $(\bm{x}_i, y_i)$, trong đó $\bm{x}_i$ là vectơ đặc trưng và $y_i$ là
nhãn đã biết. Mục tiêu là học một ánh xạ $f: \mathcal{X} \to \mathcal{Y}$ sao cho
$f(\bm{x})$ xấp xỉ tốt nhãn thật trên dữ liệu mới. Khi $\mathcal{Y}$ là tập hữu
hạn các lớp, bài toán là \textit{phân loại}; khi $\mathcal{Y}$ liên tục, bài toán
là \textit{hồi quy}.

\paragraph{Học không giám sát (\eng{unsupervised learning}).} Dữ liệu chỉ gồm
$\bm{x}_i$ mà không có nhãn. Thuật toán tìm cấu trúc tiềm ẩn trong dữ liệu: phân
cụm khách hàng thành các nhóm tương đồng, giảm chiều dữ liệu, hoặc phát hiện bất
thường. Trong lĩnh vực tín dụng, nhóm phương pháp này thường được dùng cho phân
khúc khách hàng và phát hiện gian lận.

\paragraph{Học bán giám sát (\eng{semi-supervised learning}).} Kết hợp một lượng
nhỏ dữ liệu có nhãn với lượng lớn dữ liệu không nhãn. Trong bài toán tín dụng,
tình huống này xuất hiện tự nhiên: các hồ sơ \textit{bị từ chối} không có kết quả
trả nợ nên không có nhãn, trong khi chúng chiếm một phần đáng kể tổng số hồ sơ
nộp vào. Vấn đề này được biết đến với tên gọi \eng{reject inference}.

\paragraph{Học tăng cường (\eng{reinforcement learning}).} Tác tử học chính sách
hành động thông qua tương tác với môi trường và tín hiệu thưởng/phạt. Phương pháp
này ít được dùng cho chấm điểm tín dụng do khó thiết lập môi trường thử nghiệm an
toàn, nhưng có ứng dụng trong tối ưu hạn mức và chiến lược thu hồi nợ.

Bài toán của đề án -- dự báo nhãn nhị phân ``tốt/xấu'' từ dữ liệu lịch sử đã có
kết quả -- thuộc nhóm \textbf{học có giám sát, bài toán phân loại nhị phân}. Tuy
nhiên, cần lưu ý một đặc thù: kết quả cần có không phải là nhãn dự báo
mà là \textit{xác suất} vỡ nợ được hiệu chuẩn đúng, vì xác suất này được dùng
trực tiếp để phân vùng quyết định và tính tổn thất kỳ vọng theo
công thức~\eqref{eq:expected-loss}.

\subsection{Học máy có giám sát trong bài toán tín dụng}
\label{subsec:hoc-co-giam-sat}

Phát biểu hình thức bài toán như sau. Cho tập huấn luyện
$\mathcal{D} = \{(\bm{x}_i, y_i)\}_{i=1}^{n}$ với $\bm{x}_i \in \mathbb{R}^d$ là
vectơ đặc trưng của hồ sơ thứ $i$ và $y_i \in \{0, 1\}$ là nhãn ($y_i = 1$ nếu hồ
sơ phát sinh nợ xấu). Cần tìm hàm $f$ ước lượng xác suất có điều kiện

\begin{equation}
\PD(\bm{x}) = \Pr(y = 1 \mid \bm{X} = \bm{x}).
\label{eq:pd-definition}
\end{equation}

Hàm $f$ được chọn bằng cách tối thiểu hóa một hàm mất mát trên tập huấn luyện, có
bổ sung thành phần điều chuẩn để kiểm soát độ phức tạp:

\begin{equation}
\hat{f} = \argmin_{f \in \mathcal{F}}
\left[ \frac{1}{n} \sum_{i=1}^{n} \ell\big(y_i, f(\bm{x}_i)\big)
       + \lambda \, \Omega(f) \right].
\label{eq:erm}
\end{equation}

Hàm mất mát chuẩn cho bài toán phân loại nhị phân có đầu ra xác suất là hàm
\textbf{entropy chéo nhị phân} (\eng{binary cross-entropy}, còn gọi là log-loss):

\begin{equation}
\ell(y, p) = -\big[\, y \log p + (1 - y) \log (1 - p) \,\big].
\label{eq:logloss}
\end{equation}

Log-loss thuộc lớp \textit{quy tắc chấm điểm chính đáng} (\eng{strictly proper
scoring rule}): kỳ vọng của nó đạt cực tiểu khi và chỉ khi $p$ bằng đúng xác suất
thật. Tính chất này có ý nghĩa quyết định đối với bài toán tín dụng, nơi giá trị
xác suất được sử dụng trực tiếp chứ không chỉ dùng để xếp hạng. Nếu thay bằng một
hàm mất mát không chính đáng, chẳng hạn tối ưu trực tiếp độ chính xác phân loại,
mô hình có thể xếp hạng tốt nhưng cho ra giá trị xác suất sai lệch hệ thống.

\paragraph{Ba đặc thù của bài toán tín dụng.} So với bài toán phân loại nhị phân
thông thường, bài toán chấm điểm tín dụng có ba đặc thù cần được xử lý riêng:

\begin{enumerate}[label=(\arabic*)]
  \item \textbf{Dữ liệu mất cân bằng.} Tỷ lệ hồ sơ xấu thường ở mức 3--20\%, tùy
  sản phẩm và phân khúc. Mức mất cân bằng này khiến chỉ tiêu độ chính xác
  (\eng{accuracy}) mất ý nghĩa: một mô hình dự báo ``mọi hồ sơ đều tốt'' vẫn đạt
  độ chính xác 85\% mà không có giá trị sử dụng nào.

  \item \textbf{Chi phí hai loại sai lầm rất khác nhau.} Bỏ sót một hồ sơ xấu
  (\eng{false negative}) gây tổn thất bằng một phần đáng kể dư nợ, trong khi từ
  chối sai một hồ sơ tốt (\eng{false positive}) chỉ làm mất phần lãi biên. Tỷ lệ
  chi phí thường được ước lượng trong khoảng 10:1 đến 15:1. Hệ quả là ngưỡng
  quyết định phải được chọn theo hàm chi phí~\cite{elkan2001} chứ không phải theo
  giá trị mặc định 0,5.

  \item \textbf{Trôi dạt theo thời gian.} Phân phối đặc trưng và quan hệ giữa đặc
  trưng với nhãn thay đổi theo chu kỳ kinh tế, theo chính sách sản phẩm và theo
  thành phần khách hàng~\cite{gama2014}. Do đó việc đánh giá mô hình phải được
  thực hiện trên dữ liệu của giai đoạn \textit{sau} giai đoạn huấn luyện, chứ
  không phải trên mẫu ngẫu nhiên trộn lẫn các giai đoạn.
\end{enumerate}

\subsection{Quy trình xây dựng mô hình học máy}
\label{subsec:quy-trinh-xay-dung-mo-hinh}

Quy trình xây dựng một mô hình học máy phục vụ nghiệp vụ tín dụng gồm bảy bước,
tổ chức thành ba khối như minh họa ở Hình~\ref{fig:ml-pipeline}: khối
\textit{chuẩn bị dữ liệu} (bước 1--3), khối \textit{xây dựng mô hình} (bước 4--6)
và khối \textit{vận hành} (bước 7).

\begin{figure}[htbp]
\centering
\begin{tikzpicture}[
  font=\small,
  module/.style={draw=blue!45!black, fill=blue!9, rounded corners=3pt,
                 line width=0.6pt, blur shadow={shadow blur steps=5}},
  panel/.style={draw=black!55, rounded corners=3pt, line width=0.5pt},
  inner/.style={draw=black!60, fill=white, rounded corners=2pt,
                align=center, font=\footnotesize, inner sep=2pt,
                minimum height=0.95cm, text width=2.35cm},
  flow/.style={-{Stealth[length=2.2mm]}, line width=0.8pt, draw=black!70},
  up/.style={-{Triangle[length=2.2mm, width=4.2mm]}, line width=2.6pt,
             draw=orange!75!yellow},
]
% ---------------- khối trung tâm: xây dựng mô hình -----------------------
\draw[module] (0,0) rectangle (8.6,7.2);
\node[font=\small\bfseries] at (4.3,6.88) {Khối xây dựng mô hình học máy};

% tầng trên: bước (4)
\draw[panel, fill=yellow!28] (0.3,4.55) rectangle (8.3,6.55);
\node[font=\small] at (4.3,6.25) {(4) Trích xuất và lựa chọn đặc trưng};
\node[inner] at (1.60,5.22) {Đặc trưng\\theo cửa sổ};
\node[inner] at (4.30,5.22) {Tỷ số và\\tương tác};
\node[inner] at (7.00,5.22) {Lọc IV, PSI,\\tương quan};

% tầng giữa: bước (5)
\node[draw=black!55, fill=orange!18, rounded corners=3pt, align=center,
      minimum width=7.0cm, minimum height=1.0cm] (train) at (4.3,3.55)
      {(5) Huấn luyện và tối ưu mô hình\\[-1pt]
       {\footnotesize trọng số lớp · dừng sớm · tìm kiếm ngẫu nhiên}};

% tầng dưới: bước (6)
\draw[panel, fill=green!22] (0.3,0.3) rectangle (8.3,2.55);
\node[font=\small] at (4.3,2.25) {(6) Đánh giá và hiệu chuẩn};
\node[inner] at (1.60,1.15) {AUC, Gini,\\KS, ECE};
\node[inner] at (4.30,1.15) {Hiệu chuẩn\\xác suất};
\node[inner] at (7.00,1.15) {Chọn ngưỡng\\theo chi phí};

% mũi tên giữa các tầng
\draw[up] (4.3,4.55) -- (4.3,4.09);
\draw[up] (4.3,3.03) -- (4.3,2.58);

% ---------------- khối trái: chuẩn bị dữ liệu -----------------------------
\draw[module] (-2.3,0) rectangle (-0.55,7.2);
\node[rotate=90, align=center] at (-1.425,3.6)
  {\textbf{Khối chuẩn bị dữ liệu}\\[2pt]
   {\footnotesize (1) Xác định bài toán và biến mục tiêu}\\[1pt]
   {\footnotesize (2) Thu thập \,·\, (3) Tiền xử lý, làm sạch}};

% ---------------- khối phải: vận hành -------------------------------------
\draw[module] (9.15,0) rectangle (10.9,7.2);
\node[rotate=90, align=center] at (10.025,3.6)
  {\textbf{Khối vận hành}\\[2pt]
   {\footnotesize (7) Triển khai và giám sát}\\[1pt]
   {\footnotesize dịch vụ REST \,·\, giám sát PSI}};

% ---------------- đầu vào / đầu ra ---------------------------------------
\node[rotate=90, font=\footnotesize] at (-3.0,3.6) {Dữ liệu hồ sơ tín dụng};
\node[rotate=90, font=\footnotesize] at (11.6,3.6) {PD, khuyến nghị và lý do};
\draw[flow] (-2.77,3.6) -- (-2.3,3.6);
\draw[flow] (-0.55,3.6) -- (0,3.6);
\draw[flow] (8.6,3.6) -- (9.15,3.6);
\draw[flow] (10.9,3.6) -- (11.37,3.6);

% ---------------- vòng phản hồi -------------------------------------------
\draw[flow, dashed, draw=black!55] (10.025,0) -- (10.025,-0.55)
  -- (-1.425,-0.55) -- (-1.425,0);
\node[font=\footnotesize\itshape, fill=white, inner sep=2pt] at (4.3,-0.55)
  {cập nhật dữ liệu và huấn luyện lại khi phát hiện suy giảm};
\end{tikzpicture}
\caption{Quy trình xây dựng và vận hành mô hình học máy cho bài toán tín dụng}
\label{fig:ml-pipeline}
\end{figure}

Bước (1) định nghĩa thế nào là hồ sơ ``xấu'', xác định cửa sổ quan sát và phạm vi
khách hàng áp dụng; đây là bước có ảnh hưởng lớn nhất tới toàn bộ kết quả, vì một
định nghĩa nhãn không phù hợp sẽ làm sai lệch mọi bước tiếp theo. Các bước (2)
đến (6) lần lượt tổng hợp dữ liệu về cùng một đơn vị phân tích, làm sạch, tạo đặc
trưng, huấn luyện và đánh giá mô hình. Bước (7) theo dõi mô hình trong vận hành để
xác định thời điểm huấn luyện lại. Nội dung chi tiết của từng bước, gắn với bài
toán cụ thể của đề án, được trình bày ở Chương~\ref{chap:xay-dung-he-thong}.

Quy trình này có \textit{tính lặp}: kết quả ở bước sau thường buộc phải quay lại
điều chỉnh các bước trước. Chẳng hạn, một đặc trưng có chỉ số Information Value
cao bất thường ở bước (4) có thể là dấu hiệu biến đó chứa thông tin của tương
lai, buộc phải xem lại cách tổng hợp dữ liệu ở bước (2).

\subsection{Vai trò của học máy trong hỗ trợ kiểm duyệt tín dụng}
\label{subsec:vai-tro-hoc-may}

Trong quy trình kiểm duyệt mà đề án xây dựng, mô hình học máy đảm nhiệm bốn vai
trò. Thứ
nhất, \textit{lượng hóa rủi ro trên thang đo liên tục}: thay vì kết luận nhị phân
``đạt/không đạt'', mô hình cho ra giá trị $\PD \in (0, 1)$, nhờ đó có thể phân
hạng rủi ro, định giá theo rủi ro và phân bổ nguồn lực thẩm định theo mức độ cần
thiết. Thứ hai, \textit{phát hiện quan hệ mà bộ quy tắc bỏ sót}, kể cả những
tương tác mà người xây dựng chính sách không dự liệu trước. Thứ ba, \textit{bảo
đảm tính nhất quán}: cùng một hồ sơ luôn nhận được cùng một kết quả đánh giá,
không phụ thuộc người xử lý hay thời điểm xử lý. Thứ tư, \textit{giảm khối lượng
công việc thủ công} bằng cách tự động xử lý các hồ sơ nằm ở hai vùng biên, để cán
bộ tín dụng tập trung vào nhóm hồ sơ thực sự cần cân nhắc.

Đi kèm là bốn giới hạn. Mô hình chỉ học được quy luật có trong dữ liệu lịch sử,
nên những tình huống chưa từng xuất hiện nằm ngoài khả năng dự báo. Chất lượng mô
hình phụ thuộc trực tiếp vào chất lượng dữ liệu. Các mô hình phức tạp khó giải
thích hơn thẻ điểm truyền thống, đòi hỏi bổ sung công cụ giải thích. Cuối cùng,
mô hình suy giảm theo thời gian và cần được giám sát, hiệu chỉnh định
kỳ~\cite{sculley2015}. Bốn giới hạn này là cơ sở để đề án xác định vai trò của mô
hình là \textit{hỗ trợ} chứ không phải \textit{thay thế} quá trình kiểm duyệt của
con người.

\section{Dữ liệu và đặc trưng trong bài toán tín dụng}
\label{sec:du-lieu-dac-trung}

\subsection{Đặc điểm dữ liệu tín dụng}
\label{subsec:dac-diem-du-lieu}

Dữ liệu phục vụ bài toán chấm điểm tín dụng có năm đặc điểm gây khó khăn riêng so
với dữ liệu của nhiều lĩnh vực khác.

\paragraph{Tính đúng đắn theo thời điểm.} Mọi đặc trưng của một hồ sơ chỉ được
phép tính từ dữ liệu đã tồn tại \textit{trước hoặc tại} thời điểm nộp hồ sơ. Vi
phạm nguyên tắc này dẫn tới hiện tượng rò rỉ nhãn (\eng{label leakage}): mô hình
sử dụng thông tin của tương lai, cho kết quả đánh giá cao giả tạo trong môi
trường thử nghiệm nhưng suy giảm nghiêm trọng khi triển khai. Đây là lỗi khó phát
hiện, vì nó không sinh ra thông báo lỗi mà chỉ làm các chỉ tiêu chất lượng cao
lên bất thường.

\paragraph{Độ trễ của nhãn.} Kết quả trả nợ chỉ được biết sau cửa sổ quan sát,
thường là 12 tháng. Do đó tại mọi thời điểm luôn tồn tại một khối hồ sơ gần đây
chưa có nhãn. Điều này giới hạn lượng dữ liệu huấn luyện khả dụng và buộc phải
cân nhắc đánh đổi: cửa sổ quan sát dài hơn cho nhãn đáng tin cậy hơn nhưng làm dữ
liệu huấn luyện cũ hơn.

\paragraph{Mất cân bằng lớp.} Tỷ lệ hồ sơ xấu thường thấp, khiến số quan sát của
lớp thiểu số -- vốn là lớp cần quan tâm -- bị hạn chế.

\paragraph{Giá trị khuyết mang thông tin.} Trong dữ liệu tín dụng, việc một
trường bị trống thường không phải lỗi thu thập mà tự nó là một tín hiệu. Khách
hàng không có bất kỳ thông tin nào tại CIC (\eng{thin-file}) nghĩa là chưa từng
có quan hệ tín dụng -- đây là một trạng thái có ý nghĩa rủi ro riêng, khác hẳn
với việc không truy vấn được dữ liệu. Cách xử lý đúng là phân biệt hai
trường hợp và bổ sung biến chỉ báo khuyết thay vì đơn thuần điền giá trị thống
kê.

\paragraph{Trôi dạt theo thời gian.} Cả phân phối biên của đặc trưng lẫn quan hệ
giữa đặc trưng và nhãn đều thay đổi theo thời gian, do biến động kinh tế vĩ mô,
thay đổi chính sách sản phẩm hoặc dịch chuyển cơ cấu khách hàng.

\subsection{Các nhóm dữ liệu sử dụng trong đánh giá tín dụng}
\label{subsec:nhom-du-lieu}

Dữ liệu phục vụ đánh giá tín dụng khách hàng cá nhân được phân thành bốn nhóm
nguồn, tổng hợp tại Bảng~\ref{tab:nhom-du-lieu}.

\begin{table}[htbp]
\centering
\caption{Các nhóm nguồn dữ liệu sử dụng trong đánh giá tín dụng cá nhân}
\label{tab:nhom-du-lieu}
\small
\begin{tabular}{@{}L{3.0cm}L{4.6cm}L{4.6cm}@{}}
\toprule
\textbf{Nhóm} & \textbf{Nội dung chính} & \textbf{Đặc điểm} \\
\midrule
Thông tin cá nhân và tài chính &
Tuổi, giới tính, tình trạng hôn nhân, số người phụ thuộc, trình độ, nghề nghiệp,
thâm niên, thu nhập, nơi cư trú &
Khách hàng tự khai báo; cần đối chiếu với chứng từ. Thay đổi chậm theo thời gian
\\[2pt]
Tham số khoản vay &
Sản phẩm, số tiền, kỳ hạn, lãi suất, mục đích vay, kênh tiếp nhận, tài sản bảo
đảm &
Xác định tại thời điểm nộp hồ sơ; phản ánh cả nhu cầu lẫn mức độ chấp nhận rủi ro
của khách hàng \\[2pt]
Lịch sử tín dụng và quan hệ CIC &
Điểm CIC, nhóm nợ, số ngày quá hạn, số tổ chức có quan hệ, dư nợ, tỷ lệ sử dụng
hạn mức, số lần truy vấn &
Nguồn thông tin khách quan, có năng lực dự báo cao nhất; phụ thuộc dịch vụ bên
ngoài và không có với khách hàng thin-file \\[2pt]
Hành vi giao dịch tài khoản &
Số lượng và giá trị giao dịch, dòng tiền vào/ra, tính đều đặn của thu nhập, cơ
cấu chi tiêu theo nhóm &
Chỉ có với khách hàng đã có quan hệ tài khoản; cập nhật liên tục và phản ánh tình
trạng tài chính hiện tại \\
\bottomrule
\end{tabular}
\end{table}

Trong bốn nhóm trên, nhóm dữ liệu hành vi giao dịch có ba đặc điểm riêng: tổ chức
tín dụng \textit{tự sở hữu} nên không phụ thuộc bên thứ ba, phản ánh tình trạng
tài chính \textit{hiện tại} thay vì lịch sử, và luôn có sẵn kể cả với khách hàng
thiếu lịch sử tín dụng. Onay và Öztürk~\cite{onay2018} nhấn mạnh vai trò ngày
càng lớn của nhóm dữ liệu này trong bối cảnh dữ liệu lớn, đặc biệt với phân khúc
khách hàng mà các nguồn dữ liệu truyền thống gần như không cung cấp thông tin.
Một trong những nội dung thực nghiệm của đề án là định lượng mức đóng góp của
nhóm dữ liệu này vào chất lượng dự báo.

\subsection{Xây dựng và lựa chọn đặc trưng}
\label{subsec:xay-dung-dac-trung}

\paragraph{Xây dựng đặc trưng.} Dữ liệu thô hiếm khi ở dạng phù hợp trực tiếp cho
mô hình. Quá trình xây dựng đặc trưng (\eng{feature engineering}) biến đổi dữ
liệu thô thành các biến mang ý nghĩa dự báo, thông qua bốn kỹ thuật chính:

\begin{itemize}
  \item \textbf{Tổng hợp theo cửa sổ thời gian}: từ danh sách giao dịch, tính
  tổng dòng tiền vào, số lần giao dịch, giá trị trung bình trong các cửa sổ 7,
  30, 90 và 180 ngày trước thời điểm nộp hồ sơ.

  \item \textbf{Tạo tỷ số}: các chỉ tiêu nghiệp vụ như $\DTI$, $\LTV$, tỷ lệ thu
  nhập còn lại sau trả nợ. Tỷ số có ưu điểm là chuẩn hóa được sự khác biệt về quy
  mô giữa các khách hàng.

  \item \textbf{Tạo biến xu thế}: so sánh giá trị ở cửa sổ ngắn với cửa sổ dài để
  nắm bắt chiều hướng thay đổi, chẳng hạn dòng tiền vào 30 ngày gần nhất so với
  mức trung bình 90 ngày.

  \item \textbf{Tạo biến tương tác}: kết hợp hai hoặc nhiều biến để biểu diễn
  quan hệ phi cộng tính, chẳng hạn tích của số ngày quá hạn với số tổ chức tín
  dụng có quan hệ.
\end{itemize}

\paragraph{Weight of Evidence và Information Value.} Trong ngành tín dụng, hai
công cụ được sử dụng phổ biến để đánh giá và biến đổi đặc trưng là Weight of
Evidence ($\WOE$) và Information Value ($\IV$)~\cite{siddiqi2017}. Với một biến
được rời rạc hóa thành $K$ khoảng, $\WOE$ của khoảng thứ $k$ được định nghĩa:

\begin{equation}
\WOE_k = \ln \frac{\text{tỷ lệ hồ sơ tốt trong khoảng } k}
                  {\text{tỷ lệ hồ sơ xấu trong khoảng } k}
       = \ln \frac{g_k / G}{b_k / B},
\label{eq:woe}
\end{equation}

trong đó $g_k, b_k$ lần lượt là số hồ sơ tốt và xấu trong khoảng $k$; $G, B$ là
tổng số hồ sơ tốt và xấu. Information Value tổng hợp năng lực phân biệt của toàn
bộ biến:

\begin{equation}
\IV = \sum_{k=1}^{K} \left( \frac{g_k}{G} - \frac{b_k}{B} \right) \WOE_k.
\label{eq:iv}
\end{equation}

Thực hành ngành sử dụng thang diễn giải: $\IV < 0{,}02$ -- biến không có khả năng
phân biệt; $0{,}02 \le \IV < 0{,}1$ -- yếu; $0{,}1 \le \IV < 0{,}3$ -- trung
bình; $0{,}3 \le \IV < 0{,}5$ -- mạnh; $\IV \ge 0{,}5$ -- rất mạnh, cần kiểm tra
khả năng rò rỉ nhãn.

Phép biến đổi $\WOE$ có ba ưu điểm đối với bài toán tín dụng: nó đưa mọi biến về
cùng một thang đo theo log-odds, xử lý được giá trị khuyết như một khoảng riêng,
và cho phép áp đặt ràng buộc đơn điệu để bảo đảm mô hình phù hợp với trực giác
nghiệp vụ (điểm tín dụng càng cao thì rủi ro càng thấp).

\paragraph{Lựa chọn đặc trưng.} Khi tập đặc trưng ban đầu lớn, việc giữ lại toàn
bộ gây ba vấn đề: hệ số của mô hình tuyến tính mất ổn định do đa cộng tuyến, chi
phí tính toán tăng, và khả năng diễn giải giảm. Các phương pháp lựa chọn đặc
trưng được chia thành ba nhóm:

\begin{itemize}
  \item \textbf{Phương pháp lọc} (\eng{filter}): đánh giá từng đặc trưng độc lập
  với mô hình, dựa trên $\IV$, hệ số tương quan hoặc kiểm định thống kê. Chi phí
  tính toán thấp, phù hợp để sàng lọc sơ bộ.

  \item \textbf{Phương pháp bao} (\eng{wrapper}): thử các tập con đặc trưng và
  đánh giá bằng chính mô hình đích. Cho kết quả tốt hơn nhưng chi phí tính toán
  cao.

  \item \textbf{Phương pháp nhúng} (\eng{embedded}): việc lựa chọn diễn ra trong
  quá trình huấn luyện, chẳng hạn điều chuẩn $L_1$ hoặc tầm quan trọng đặc trưng
  của mô hình cây.
\end{itemize}

Đề án sử dụng quy trình lọc ba bước nối tiếp -- theo năng lực dự báo ($\IV$),
theo mức tương quan, và theo độ ổn định theo thời gian ($\PSI$) -- được trình bày
chi tiết ở mục~\ref{sec:trich-xuat-dac-trung}.

\section{Các mô hình học máy áp dụng trong đánh giá tín dụng}
\label{sec:cac-mo-hinh}

Mục này trình bày nguyên lý của năm nhóm thuật toán được đề án nghiên cứu và thực
nghiệm, theo đúng thứ tự từ đơn giản đến phức tạp.

\subsection{Hồi quy Logistic}
\label{subsec:hoi-quy-logistic}

Hồi quy logistic là mô hình tuyến tính tổng quát dùng cho biến phụ thuộc nhị
phân. Mô hình giả định log-odds của xác suất vỡ nợ là hàm tuyến tính của các đặc
trưng:

\begin{equation}
\ln \frac{\PD(\bm{x})}{1 - \PD(\bm{x})} = \beta_0 + \sum_{j=1}^{d} \beta_j x_j,
\qquad\text{tương đương}\qquad
\PD(\bm{x}) = \frac{1}{1 + e^{-(\beta_0 + \bm{\beta}^{\!\top}\bm{x})}}.
\label{eq:logistic}
\end{equation}
\vspace{-0.6\baselineskip}

Tham số $\bm{\beta}$ được ước lượng bằng phương pháp hợp lý cực đại, tương đương
với việc tối thiểu hóa hàm log-loss~\eqref{eq:logloss}. Để kiểm soát quá khớp,
thường bổ sung thành phần điều chuẩn $L_2$:

\begin{equation}
\hat{\bm{\beta}} = \argmin_{\bm{\beta}}
\left[ -\sum_{i=1}^{n} \Big( y_i \log p_i + (1-y_i)\log(1-p_i) \Big)
       + \frac{1}{2C} \lVert \bm{\beta} \rVert_2^2 \right].
\label{eq:logistic-l2}
\end{equation}

\paragraph{Ưu điểm.} Hồi quy logistic có ba ưu thế khiến nó vẫn là mô hình chuẩn
mực trong ngành tín dụng sau hơn nửa thế kỷ. Thứ nhất, \textit{khả năng giải
thích}: mỗi hệ số $\beta_j$ cho biết trực tiếp mức thay đổi của log-odds khi đặc
trưng $x_j$ tăng một đơn vị, và dấu của hệ số có thể được đối chiếu với kỳ vọng
nghiệp vụ. Thứ hai, \textit{xác suất đầu ra được hiệu chuẩn tốt} do hàm mất mát
chính là log-loss. Thứ ba, \textit{tính ổn định}: mô hình có ít tham số nên ít
nhạy cảm với nhiễu trong dữ liệu huấn luyện.

\paragraph{Hạn chế.} Giả định tuyến tính trên thang log-odds là hạn chế cơ bản:
mô hình không tự nắm bắt được quan hệ phi tuyến hay tương tác giữa các biến, trừ
khi những quan hệ này được đưa vào tường minh dưới dạng biến phái sinh.

\paragraph{Thẻ điểm dựa trên biến đổi WOE.} Trong thực tiễn ngân hàng, hồi quy
logistic thường được áp dụng trên các đặc trưng đã qua biến đổi
$\WOE$~\eqref{eq:woe} thay vì trên giá trị gốc. Cách làm này khắc phục một phần
hạn chế tuyến tính -- vì việc rời rạc hóa cho phép biểu diễn quan hệ phi tuyến
theo từng khoảng -- đồng thời cho phép quy đổi kết quả thành thang điểm nguyên
dễ vận hành:

\begin{equation}
\text{Điểm} = \text{Offset} + \text{Factor} \times \ln(\text{odds}),
\qquad
\text{Factor} = \frac{\text{PDO}}{\ln 2},
\label{eq:scorecard}
\end{equation}

trong đó PDO (\eng{Points to Double the Odds}) là số điểm tương ứng với việc tỷ
số odds tăng gấp đôi. Thẻ điểm xây dựng theo cách này có thể in ra thành bảng
tra, giúp cán bộ tín dụng hiểu chính xác vì sao một hồ sơ đạt mức điểm cụ thể.

\subsection{Decision Tree}
\label{subsec:decision-tree}

Cây quyết định~\cite{breiman1984} phân chia đệ quy không gian đặc trưng thành các
vùng con, tại mỗi
bước chọn cặp (đặc trưng, ngưỡng) làm giảm nhiều nhất độ không thuần khiết của
nút. Với bài toán phân loại, độ không thuần khiết thường được đo bằng chỉ số
Gini:

\begin{equation}
G(t) = 1 - \sum_{c \in \{0,1\}} p_c(t)^2,
\label{eq:gini-impurity}
\end{equation}

trong đó $p_c(t)$ là tỷ lệ quan sát thuộc lớp $c$ tại nút $t$. Phép chia tốt nhất
là phép chia làm cực đại mức giảm độ không thuần khiết có trọng số:

\begin{equation}
\Delta G = G(t) - \frac{n_L}{n_t} G(t_L) - \frac{n_R}{n_t} G(t_R).
\label{eq:gini-gain}
\end{equation}

\paragraph{Ưu điểm.} Cây quyết định tự động nắm bắt quan hệ phi tuyến và tương
tác giữa các biến, không yêu cầu chuẩn hóa thang đo, xử lý được cả biến định
lượng lẫn định tính, và -- với cây nông -- có thể diễn giải trực tiếp thành chuỗi
điều kiện gần với ngôn ngữ nghiệp vụ.

\paragraph{Hạn chế.} Cây đơn có phương sai cao: một thay đổi nhỏ trong dữ liệu
huấn luyện có thể dẫn tới cấu trúc cây hoàn toàn khác. Cây sâu dễ quá khớp, trong
khi cây nông lại quá đơn giản để nắm bắt quan hệ phức tạp. Để kiểm soát, đề án sử
dụng kỹ thuật tỉa nhánh theo độ phức tạp chi phí (\eng{cost-complexity pruning})
với tiêu chí:

\begin{equation}
R_\alpha(T) = R(T) + \alpha \, |\tilde{T}|,
\label{eq:ccp}
\end{equation}

trong đó $R(T)$ là sai số của cây $T$, $|\tilde{T}|$ là số nút lá và $\alpha$ là
tham số phạt độ phức tạp.

Chính hạn chế về phương sai cao của cây đơn là động lực dẫn tới hai nhóm phương
pháp học tập hợp trình bày ở hai mục tiếp theo.

\subsection{Random Forest}
\label{subsec:random-forest}

Rừng ngẫu nhiên~\cite{breiman2001} xây dựng nhiều cây quyết định độc lập và tổng
hợp kết quả bằng cách lấy trung bình xác suất dự báo:

\begin{equation}
\PD_{\text{RF}}(\bm{x}) = \frac{1}{M} \sum_{m=1}^{M} \PD_m(\bm{x}).
\label{eq:rf}
\end{equation}

Tính đa dạng giữa các cây -- yếu tố quyết định hiệu quả của phương pháp -- được
tạo ra bằng hai cơ chế ngẫu nhiên hóa. Thứ nhất, mỗi cây được huấn luyện trên một
mẫu bootstrap lấy có hoàn lại từ tập huấn luyện. Thứ hai, tại mỗi nút, chỉ một
tập con ngẫu nhiên gồm $\sqrt{d}$ đặc trưng được xem xét để tìm phép chia.

Cơ sở lý thuyết của phương pháp~\cite{hastie2009} nằm ở công thức phương sai của trung bình $M$
biến ngẫu nhiên có cùng phương sai $\sigma^2$ và tương quan cặp $\rho$:

\begin{equation}
\operatorname{Var}\!\left( \frac{1}{M}\sum_{m=1}^{M} Z_m \right)
= \rho \sigma^2 + \frac{1 - \rho}{M}\sigma^2.
\label{eq:rf-variance}
\end{equation}

Số hạng thứ hai giảm về 0 khi $M$ tăng, nhưng số hạng thứ nhất chỉ giảm được bằng
cách làm giảm tương quan $\rho$ giữa các cây. Đây chính là lý do của cơ chế lấy
mẫu ngẫu nhiên đặc trưng tại mỗi nút: nếu mọi cây đều được phép dùng đặc trưng
mạnh nhất, chúng sẽ giống nhau và việc lấy trung bình mất tác dụng.

\paragraph{Ưu điểm.} Rừng ngẫu nhiên giữ được khả năng biểu diễn phi tuyến của
cây nhưng có phương sai thấp hơn nhiều, ít nhạy cảm với siêu tham số, và cung cấp
sẵn thước đo tầm quan trọng của đặc trưng.

\paragraph{Hạn chế.} Mô hình mất khả năng diễn giải trực tiếp; chi phí bộ nhớ và
thời gian suy luận tăng theo số cây; và xác suất đầu ra thường bị lệch về phía
trung tâm, cần hiệu chuẩn lại trước khi dùng làm $\PD$.

\subsection{Support Vector Machine}
\label{subsec:svm}

Máy vectơ hỗ trợ (SVM)~\cite{cortes1995,vapnik1998} tiếp cận bài toán phân loại
từ góc nhìn hình học: tìm siêu
phẳng phân tách hai lớp sao cho \textit{lề} -- khoảng cách từ siêu phẳng tới các
điểm gần nhất của mỗi lớp -- là lớn nhất. Với dữ liệu không phân tách tuyến tính
hoàn hảo, bài toán được phát biểu ở dạng lề mềm:

\begin{equation}
\begin{aligned}
&\min_{\bm{w}, b, \bm{\xi}} \quad
  \frac{1}{2}\lVert \bm{w} \rVert^2 + C \sum_{i=1}^{n} \xi_i \\[2pt]
&\text{với ràng buộc} \quad
  y_i(\bm{w}^{\!\top}\bm{x}_i + b) \ge 1 - \xi_i, \quad \xi_i \ge 0,
\end{aligned}
\label{eq:svm-primal}
\end{equation}

trong đó $\xi_i$ là biến bù cho phép một số điểm nằm sai phía, còn $C$ điều tiết
đánh đổi giữa độ rộng lề và mức chấp nhận vi phạm.

\paragraph{Thủ thuật hạt nhân.} Điểm mạnh đặc trưng của SVM nằm ở chỗ bài toán
đối ngẫu chỉ phụ thuộc vào tích vô hướng giữa các cặp quan sát. Nhờ đó có thể
thay tích vô hướng bằng một hàm hạt nhân $K(\bm{x}_i, \bm{x}_j)$ tương ứng với
tích vô hướng trong một không gian đặc trưng có số chiều cao hơn, mà không cần
tính toán tường minh phép ánh xạ. Hạt nhân được dùng phổ biến nhất là hạt nhân
Gauss (RBF):

\begin{equation}
K(\bm{x}_i, \bm{x}_j) = \exp\!\big( -\gamma \lVert \bm{x}_i - \bm{x}_j \rVert^2 \big).
\label{eq:rbf}
\end{equation}

\paragraph{Ưu điểm.} SVM có nền tảng lý thuyết chặt chẽ dựa trên nguyên lý cực
tiểu hóa rủi ro cấu trúc, hoạt động tốt trong không gian nhiều chiều, và nhờ thủ
thuật hạt nhân có thể biểu diễn ranh giới quyết định rất linh hoạt.

\paragraph{Hạn chế trong bài toán tín dụng.} SVM gặp ba trở ngại thực tế đáng kể
khi áp dụng cho chấm điểm tín dụng ở quy mô lớn:

\begin{enumerate}[label=(\arabic*)]
  \item \textbf{Chi phí tính toán.} Việc giải bài toán đối ngẫu có độ phức tạp từ
  $O(n^2)$ đến $O(n^3)$ theo số quan sát, khiến SVM hạt nhân đầy đủ không khả thi
  với tập dữ liệu hàng chục nghìn hồ sơ trở lên.

  \item \textbf{Đầu ra không phải xác suất.} SVM cho ra giá trị hàm quyết định --
  khoảng cách có dấu tới siêu phẳng -- chứ không phải xác suất. Để dùng được
  trong hệ thống tín dụng, giá trị này phải được quy đổi bằng phương pháp
  Platt~\cite{platt1999}: khớp một hồi quy logistic một biến trên giá trị hàm
  quyết định.

  \item \textbf{Khó giải thích.} Ranh giới quyết định trong không gian hạt nhân
  không có biểu diễn tường minh theo các đặc trưng gốc.
\end{enumerate}

Để khắc phục trở ngại thứ nhất, đề án sử dụng phép xấp xỉ
Nystroem~\cite{williams2001}: thay vì tính
ma trận hạt nhân đầy đủ kích thước $n \times n$, phương pháp chọn $m \ll n$ điểm
mốc và xây dựng một ánh xạ tường minh $\phi: \mathbb{R}^d \to \mathbb{R}^m$ sao
cho $\phi(\bm{x}_i)^{\!\top}\phi(\bm{x}_j) \approx K(\bm{x}_i, \bm{x}_j)$. Sau
phép ánh xạ này, một SVM tuyến tính -- có chi phí tuyến tính theo $n$ -- cho
nghiệm xấp xỉ nghiệm của SVM hạt nhân đầy đủ.

\subsection{Gradient Boosting}
\label{subsec:gradient-boosting}

Khác với rừng ngẫu nhiên xây dựng các cây \textit{song song và độc lập}, gradient
boosting~\cite{friedman2001} xây dựng các cây \textit{tuần tự}, trong đó mỗi cây
mới học phần sai số còn lại của tổ hợp các cây trước đó. Mô hình được xây dựng
theo sơ đồ cộng tính:

\begin{equation}
F_m(\bm{x}) = F_{m-1}(\bm{x}) + \eta \, h_m(\bm{x}),
\label{eq:boosting}
\end{equation}

trong đó $h_m$ là cây thứ $m$ được khớp với gradient âm của hàm mất mát tại dự
báo hiện tại, còn $\eta \in (0, 1]$ là tốc độ học đóng vai trò điều chuẩn: giá
trị $\eta$ nhỏ buộc mô hình cải thiện từ từ, giảm nguy cơ quá khớp nhưng đòi hỏi
nhiều cây hơn.

\paragraph{Các biến thể hiện đại.} Ba thư viện được sử dụng rộng rãi hiện nay cải
tiến ý tưởng gốc theo những hướng khác nhau:

\begin{itemize}
  \item \textbf{XGBoost}~\cite{chen2016} bổ sung thành phần điều chuẩn tường minh
  vào hàm mục tiêu và sử dụng khai triển Taylor bậc hai để xác định phép chia tối
  ưu, thay vì chỉ dùng gradient bậc nhất.

  \item \textbf{LightGBM}~\cite{ke2017} sinh cây theo chiều lá
  (\eng{leaf-wise}) -- mở rộng lá có mức giảm mất mát lớn nhất -- thay vì theo
  chiều sâu, kết hợp với kỹ thuật rời rạc hóa đặc trưng thành histogram để tăng
  tốc độ huấn luyện.

  \item \textbf{CatBoost}~\cite{prokhorenkova2018} xử lý đặc trưng định tính bằng
  thống kê mục tiêu có thứ tự (\eng{ordered target statistics}), khắc phục hiện
  tượng rò rỉ mục tiêu vốn xảy ra khi mã hóa biến định tính theo giá trị trung
  bình của nhãn.
\end{itemize}

\paragraph{Ưu điểm.} Gradient boosting thường cho chất lượng dự báo cao nhất
trong các bài toán dữ liệu dạng bảng, nhờ khả năng nắm bắt tương tác bậc cao và
cơ chế tập trung vào các quan sát khó phân loại.

\paragraph{Hạn chế.} Mô hình nhạy cảm với siêu tham số hơn rừng ngẫu nhiên, dễ
quá khớp nếu không dừng sớm, khó diễn giải trực tiếp, và xác suất đầu ra cần được
hiệu chuẩn -- đặc biệt khi sử dụng trọng số lớp để bù mất cân bằng dữ liệu.

\subsection{So sánh và lựa chọn mô hình phù hợp}
\label{subsec:so-sanh-mo-hinh}

\subsubsection{Hệ thống chỉ tiêu đánh giá}

Việc so sánh các mô hình đòi hỏi một bộ chỉ tiêu phản ánh đầy đủ các khía
cạnh mà bài toán quan tâm. Đề án sử dụng bốn nhóm chỉ tiêu.

\paragraph{Nhóm 1 -- Chỉ tiêu phân loại phổ quát.} Xuất phát từ ma trận nhầm lẫn
với bốn ô TP, FP, TN, FN:

\begin{equation}
\text{Precision} = \frac{TP}{TP + FP}, \qquad
\text{Recall} = \frac{TP}{TP + FN},
\label{eq:prf}
\end{equation}
\begin{equation}
F_1 = \frac{2 \cdot \text{Precision} \cdot \text{Recall}}
           {\text{Precision} + \text{Recall}}.
\label{eq:f1}
\end{equation}

Độ chính xác (\eng{accuracy}) được báo cáo nhưng không dùng làm tiêu chí chính,
vì lý do đã nêu ở mục~\ref{subsec:hoc-co-giam-sat}.

\paragraph{Nhóm 2 -- Chỉ tiêu đặc thù ngành tín dụng.} Ba chỉ tiêu dưới đây được
sử dụng phổ biến trong kiểm định mô hình tín dụng~\cite{baesens2016}:

\begin{itemize}
  \item \textbf{$\AUC$} -- diện tích dưới đường cong ROC, bằng xác suất một hồ sơ
  xấu chọn ngẫu nhiên nhận điểm rủi ro cao hơn một hồ sơ tốt chọn ngẫu nhiên. Chỉ
  tiêu này không phụ thuộc ngưỡng quyết định, nên phù hợp để so sánh mô hình.

  \item \textbf{Hệ số Gini} $= 2 \cdot \AUC - 1$, là cách biểu diễn quen thuộc
  hơn trong ngành ngân hàng, với thang giá trị từ 0 (không phân biệt) đến 1 (phân
  biệt hoàn hảo).

  \item \textbf{Thống kê $\KS$} -- khoảng cách lớn nhất giữa hàm phân phối tích
  lũy của điểm số ở nhóm tốt và nhóm xấu:
  \begin{equation}
  \KS = \max_{s} \big| F_{\text{tốt}}(s) - F_{\text{xấu}}(s) \big|.
  \label{eq:ks}
  \end{equation}
  Ý nghĩa vận hành của $\KS$ là nó chỉ ra điểm cắt tại đó mô hình phân tách hai
  nhóm mạnh nhất.
\end{itemize}

\paragraph{Nhóm 3 -- Chất lượng hiệu chuẩn xác suất.} Một mô hình có thể xếp hạng
rất tốt nhưng cho giá trị xác suất sai lệch. Vì đề án sử dụng $\PD$ trực tiếp
để phân vùng quyết định, chất lượng hiệu chuẩn phải được đo tường minh bằng điểm
Brier và sai số hiệu chuẩn kỳ vọng:

\begin{equation}
\text{Brier} = \frac{1}{n}\sum_{i=1}^{n} (p_i - y_i)^2,
\qquad
\text{ECE} = \sum_{b=1}^{B} \frac{n_b}{n}
             \big| \bar{y}_b - \bar{p}_b \big|.
\label{eq:brier-ece}
\end{equation}

Hai kỹ thuật hiệu chuẩn phổ biến là hồi quy đẳng hướng (\eng{isotonic
regression})~\cite{zadrozny2002} và phương pháp Platt. Nghiên cứu của
Niculescu-Mizil và Caruana~\cite{niculescu2005} cho thấy các mô hình học tập hợp
thường cho xác suất bị nén về phía trung tâm và cần hiệu chuẩn lại.

\paragraph{Nhóm 4 -- Chỉ tiêu vận hành.} Bao gồm thời gian huấn luyện, thời gian
chấm điểm một hồ sơ và mức độ giải thích được. Đây là những chỉ tiêu quyết định
khả năng triển khai thực tế nhưng thường bị bỏ qua trong các so sánh học thuật.

\subsubsection{Cơ sở lựa chọn mô hình}

Bảng~\ref{tab:so-sanh-mo-hinh} tổng hợp đặc điểm của năm nhóm mô hình theo các
tiêu chí quan trọng đối với bài toán kiểm duyệt tín dụng.

\begin{table}[htbp]
\centering
\caption{So sánh định tính các nhóm mô hình theo tiêu chí của bài toán tín dụng}
\label{tab:so-sanh-mo-hinh}
\small
\begin{tabular}{@{}L{2.6cm}ccccc@{}}
\toprule
\textbf{Mô hình} & \textbf{Năng lực} & \textbf{Giải} & \textbf{Hiệu} &
\textbf{Chi phí} & \textbf{Độ ổn} \\
 & \textbf{dự báo} & \textbf{thích} & \textbf{chuẩn} & \textbf{huấn luyện} &
 \textbf{định} \\
\midrule
Hồi quy logistic      & Trung bình & Rất cao    & Rất tốt    & Rất thấp  & Cao \\
Thẻ điểm WOE          & Trung bình & Rất cao    & Rất tốt    & Thấp      & Cao \\
Cây quyết định        & Thấp       & Cao        & Trung bình & Rất thấp  & Thấp \\
Rừng ngẫu nhiên       & Cao        & Trung bình & Cần hiệu chuẩn & Trung bình & Cao \\
Máy vectơ hỗ trợ      & Cao        & Thấp       & Cần hiệu chuẩn & Cao       & Trung bình \\
Gradient boosting     & Rất cao    & Thấp       & Cần hiệu chuẩn & Trung bình & Trung bình \\
\bottomrule
\end{tabular}
\end{table}

Bảng trên cho thấy không tồn tại một mô hình tối ưu trên mọi tiêu chí: các mô
hình có năng lực dự báo cao nhất lại là những mô hình khó giải thích nhất và cần
hiệu chuẩn xác suất. Do vậy việc lựa chọn phải dựa trên đánh đổi phù hợp với bối
cảnh sử dụng cụ thể.

Kết quả đối sánh quy mô lớn của Lessmann và cộng sự~\cite{lessmann2015} cung cấp
căn cứ định lượng cho đánh đổi này: trên nhiều bộ dữ liệu tín dụng, khoảng cách
$\AUC$ giữa mô hình học tập hợp tốt nhất và hồi quy logistic thường chỉ ở mức
0,01--0,03. Mức chênh lệch khiêm tốn này đặt ra câu hỏi liệu việc chấp nhận mất
khả năng giải thích và tăng chi phí vận hành có tương xứng với mức cải thiện đó
hay không. Câu trả lời phụ thuộc quy mô danh mục: với danh mục lớn, mức cải thiện
0,01 $\AUC$ có thể tương ứng khoản tiết kiệm đáng kể về tổn thất tín dụng; với
danh mục nhỏ, yêu cầu về tính minh bạch và chi phí vận hành thường được ưu tiên
hơn.

Trên cơ sở phân tích đó, đề án xác định nguyên tắc lựa chọn gồm bốn tiêu chí,
theo thứ tự ưu tiên: (i) năng lực phân biệt đo bằng $\AUC$ và $\KS$ trên tập kiểm
thử ngoài thời gian; (ii) chất lượng hiệu chuẩn xác suất sau khi đã hiệu chuẩn;
(iii) thời gian chấm điểm một hồ sơ trong điều kiện xử lý trực tuyến; và (iv)
khả năng giải thích kết quả cho từng hồ sơ. Việc so sánh thực nghiệm theo bốn
tiêu chí này được trình bày ở Chương~\ref{chap:thuc-nghiem}.

\subsubsection{Xử lý mất cân bằng dữ liệu}

Ba nhóm kỹ thuật thường được dùng để xử lý mất cân bằng lớp: \textit{lấy mẫu lại}
-- tăng mẫu lớp thiểu số (chẳng hạn SMOTE~\cite{chawla2002}) hoặc giảm mẫu lớp đa
số; \textit{trọng số lớp} -- gán trọng số lớn hơn cho quan sát thuộc lớp thiểu số
trong hàm mất mát; và \textit{học nhạy cảm chi phí} -- đưa ma trận chi phí vào
quá trình huấn luyện hoặc vào việc chọn ngưỡng~\cite{elkan2001}.

Đề án sử dụng \textbf{trọng số lớp} kết hợp \textbf{lựa chọn ngưỡng theo hàm chi
phí}, không sử dụng kỹ thuật lấy mẫu lại. Ở mức mất cân bằng vừa phải (khoảng
15\%), việc lấy mẫu lại làm thay đổi phân phối tiên nghiệm của dữ liệu, khiến xác
suất đầu ra bị lệch hệ thống và phải điều chỉnh lại, trong khi lợi ích về năng
lực phân biệt không rõ ràng. Trọng số lớp giữ nguyên phân phối gốc, còn phần lệch
xác suất do trọng số gây ra được xử lý ở bước hiệu chuẩn.

\subsubsection{Giải thích mô hình}

Với mô hình sử dụng trong hoạt động cấp tín dụng, khả năng giải thích là điều
kiện bắt buộc: cán bộ tín dụng cần hiểu căn cứ của khuyến nghị để quyết định chấp
nhận hay bác bỏ, còn khách hàng có quyền được biết lý do hồ sơ bị từ chối.

Phương pháp SHAP (\eng{SHapley Additive
exPlanations})~\cite{lundberg2017,lundberg2020} dựa trên giá trị Shapley trong lý
thuyết trò chơi hợp tác. Với mỗi hồ sơ, SHAP phân rã chênh lệch giữa dự báo của
mô hình và giá trị dự báo trung bình thành tổng đóng góp của từng đặc trưng:

\begin{equation}
f(\bm{x}) = \phi_0 + \sum_{j=1}^{d} \phi_j,
\label{eq:shap}
\end{equation}

trong đó $\phi_0$ là giá trị kỳ vọng của dự báo và $\phi_j$ là đóng góp của đặc
trưng thứ $j$. Tính cộng tính này cho phép chuyển kết quả SHAP thành danh sách lý
do (\eng{reason codes}) trình bày được cho người sử dụng.

\subsubsection{Giám sát trôi dạt}

Sau khi triển khai, mô hình cần được giám sát liên tục. Chỉ số được dùng phổ biến
nhất là Population Stability Index, đo mức chênh lệch phân phối của một biến giữa
giai đoạn tham chiếu và giai đoạn hiện tại:

\begin{equation}
\PSI = \sum_{k=1}^{K}
\big( p_k^{\text{hiện tại}} - p_k^{\text{tham chiếu}} \big)
\ln \frac{p_k^{\text{hiện tại}}}{p_k^{\text{tham chiếu}}}.
\label{eq:psi}
\end{equation}

Thực hành ngành sử dụng ngưỡng: $\PSI < 0{,}10$ -- phân phối ổn định;
$0{,}10 \le \PSI < 0{,}25$ -- cần theo dõi; $\PSI \ge 0{,}25$ -- thay đổi đáng
kể, cần xem xét huấn luyện lại mô hình.

\section*{Kết luận chương 1}
\addcontentsline{toc}{section}{Kết luận chương 1}

Chương 1 đã hệ thống hóa cơ sở lý luận làm nền tảng cho đề án, với bốn kết quả
chính.

Thứ nhất, chương phân biệt rành mạch giữa kiểm duyệt (quá trình đánh giá, sàng
lọc hồ sơ) với phê duyệt (hành vi ra quyết định của cấp có thẩm quyền). Sự phân
biệt này xác định vai trò \textit{hỗ trợ} của mô hình và là cơ sở cho thiết kế
ba vùng quyết định ở chương sau. Chương cũng chỉ ra bốn hạn chế mang tính cấu
trúc của phương pháp kiểm duyệt truyền thống: không biểu diễn được tương tác giữa
các yếu tố, ngưỡng đặt theo phán đoán, kết quả nhị phân thiếu thang đo mức độ, và
giới hạn về tính nhất quán cùng tốc độ xử lý.

Thứ hai, trong phân rã $\EL = \PD \times \LGD \times \EAD$, đề án xác định phạm
vi là ước lượng thành phần $\PD$, với định nghĩa vỡ nợ thống nhất theo quy định
phân loại nợ hiện hành: quá hạn từ 90 ngày trở lên, tương ứng nợ nhóm 3--5.

Thứ ba, chương phân tích ba đặc thù của bài toán tín dụng (dữ liệu mất cân bằng,
chi phí hai loại sai lầm chênh lệch lớn, trôi dạt theo thời gian) và năm đặc điểm
của dữ liệu tín dụng, trong đó quan trọng nhất là yêu cầu về tính đúng đắn theo
thời điểm. Những đặc thù này quy định ba lựa chọn phương pháp ở các chương sau:
dùng trọng số lớp thay vì lấy mẫu lại, chọn ngưỡng theo hàm chi phí thay vì giá
trị mặc định, và chia dữ liệu theo thời gian thay vì chia ngẫu nhiên.

Thứ tư, chương trình bày nguyên lý, ưu điểm và hạn chế của năm nhóm thuật toán
được nghiên cứu cùng bốn nhóm chỉ tiêu đánh giá. Phân tích cho thấy
không tồn tại mô hình tối ưu trên mọi tiêu chí, nên việc lựa chọn phải dựa trên
đánh đổi giữa năng lực dự báo, khả năng giải thích, chất lượng hiệu chuẩn và chi
phí vận hành.
